\section{本讲小结}

事件的位置和时刻必须由同一参考系的尺与同步钟定义。洛伦兹变换保持时空间隔不变，同时改变空间和时间坐标；收到光信号的时刻还包含传播延迟。类时关联的因果先后不因换系颠倒，类空事件的同时性和先后则可以依赖参考系。

时间膨胀比较同一只钟经历的两个事件，长度收缩要求在测量系同时记录两端。换系后，为测量长度通常须重选端点事件。速度变换把同一粒子相邻事件的空间差与时间差同时变换，因此不能直接相加高速速度，也不能把横向坐标不变误解为横向速度不变。

有质量粒子的总能量和动量满足 $E^2-p^2c^2=m_0^2c^4$。做功决定动能增量，相对论关系决定由动能反求的速度；恒力使动量线性增长，速度仍小于光速。碰撞、衰变与散射须在同一参考系内列出全部参与者的能量和动量，再由总量求系统不变质量。系统质量包含内部能量，通常不能逐个相加粒子的静质量。

光学实验中的条纹比较仍由往返传播时间与相位决定，介质中的光速须先注明介质静止系。局部加速效应与引力效应的等价有空间和时间范围限制；弱静态引力场中，时钟比较取决于引力势差。运动学效应与引力效应使用的条件不同，具体比较应把所选时钟、事件和测量量一并注明。
